%% --- Replace the opening paragraph with this ---

Modeling ASO binding to mRNA required both a physically stable representation
of the target RNA and a simulation framework in which ASOs could interact
with that target. This chapter describes the computational pipeline used in
this work. Experimental RNA structures and NuFold-predicted RNA structures
were first prepared. Coarse-grained molecular dynamics (CGMD) was then used
to fold and equilibrate each sequence, and candidate ASO-accessible regions
were identified using secondary-structure and solvent-accessible surface area
(SASA) analyses. ASO-RNA binding was subsequently quantified using a
center-of-mass (COM) distance criterion. All simulations were run in
LAMMPS~\cite{Thompson2022}. RNA folding and binding were modeled with the
SIS coarse-grained RNA force field as implemented by Aierken and
Joseph~\cite{Aierken2024}. Post-processing was performed with custom Python
analysis scripts.
% TODO: add the original Nguyen et al. SIS-model citation here, as requested by reviewer.

%% --- Replace 3NJ6 subsection with this ---

\subsection{3NJ6 Reference Structure}
\label{sec:3nj6}

Validation of the simulation parameters was carried out using the crystal
structure of the CAG-repeat RNA duplex 3NJ6~\cite{Aierken2024}. Because this
system has already been used to validate the SIS model parameters, it served
as a benchmark for confirming correct force-field implementation in the
present workflow. The C3$^\prime$ coordinates from the selected strand were
extracted from the experimental structure and used as the fixed reference for
root-mean-square deviation (RMSD) calculations.

%% --- Replace NuFold subsection with this ---
%% This version avoids overclaiming "Figure 5" until you verify the exact set.

\subsection{NuFold Benchmark Structures}
\label{sec:nufold}

In addition to the 3NJ6 benchmark, we analyzed a selected set of benchmark
targets reported by Kagaya et al.~\cite{Kagaya2025}. The seven targets used
here were 2DER (chain C, 76 nt), 6CK4 (chain A, 117 nt), 4QLN (rp12, 125 nt),
5ML7 (chain A, 96 nt), 4OJI (chain A, 54 nt), 6G7Z, and 4U7U (chain L,
61 nt). For each target sequence, a multiple-sequence alignment was generated
with rMSA, and a secondary-structure constraint was generated with
IPknot before running NuFold with the training-configuration preset. The
rank-1 output PDB for each target was taken as the predicted three-dimensional
structure. Each rank-1 PDB was then parsed to extract the corresponding
C3$^\prime$ coordinates, which were stored as the fixed RMSD reference for all
subsequent simulation runs on that RNA target.
% TODO: add citations for rMSA and IPknot in the bibliography.
% TODO: verify whether 6G7Z belongs in this exact benchmark set and whether
% "Figure 5" should be cited here, since the reviewer flagged this.

%% --- Replace Longer and mRNA-Scale subsection with this ---

\subsection{Longer and mRNA-Scale Sequences}
\label{sec:longer_seqs}

Ten RNA sequences with nucleotide compositions chosen to resemble naturally
occurring mRNA were randomly generated with lengths between 50 and 200
nucleotides. In addition, sequences of approximately 1000 nucleotides were
used to test whether the pipeline scaled to mRNA-length chains. All sequences
in this length range were analyzed using the simulation and secondary-structure
workflow described in Sections~\ref{sec:sim_protocol} and~\ref{sec:sec_struct},
respectively.

%% --- Replace Sequence Extraction subsection with this ---

\subsubsection{Sequence Extraction from PDB Files}

For each RNA target, the source PDB or PDB1 file was parsed to extract residue
information. ATOM and HETATM records were filtered to retain recognized
nucleic-acid residues. Residue names labeled as thymine in the input records
were mapped to uracil in the parsed sequence so that downstream topology
generation used a uniform RNA alphabet; this was a parser-level naming
standardization step rather than a chemical modification of the structure.
When multiple chains were present in a PDB file, the chain containing the
largest number of RNA residues was selected automatically. For the selected
NuFold benchmark targets, chain assignments were fixed to match those reported
by Kagaya et al.~\cite{Kagaya2025} when applicable (e.g., chain C for 2DER
and chain A for 6CK4). Residues were then sorted in ascending sequence-number
order, and the resulting nucleotide sequence (A, C, G, U) was used as the
basis for LAMMPS data-file generation.

%% --- Replace LAMMPS Data File Generation subsection with this ---

\subsubsection{LAMMPS Data File Generation}

The initial LAMMPS configuration placed coarse-grained beads linearly along
the $x$-axis with 5.9~\AA\ spacing between adjacent beads, equal to the
harmonic-bond equilibrium length. This produced a fully extended initial chain
independent of the geometry of the source PDB file. Sequential 1-2 bonds and
1-2-3 angles were created according to residue order. Three-dimensional
periodic boundary conditions were applied in a cubic simulation box with
user-specified edge length. Atom types, masses, initial coordinates, bonded
connectivity, angular connectivity, and box dimensions were written to the
LAMMPS data file.

The C3$^\prime$ coordinates extracted from the source PDB were used only as the
reference structure for RMSD calculations; they were not used as the initial
bead coordinates in the coarse-grained LAMMPS configuration.

%% --- Replace force-field opening paragraph with this ---

\section{Coarse-Grained Force Field}
\label{sec:forcefield}

All simulations used the SIS coarse-grained RNA force field as implemented by
Aierken and Joseph~\cite{Aierken2024}. In this model, each nucleotide is
represented by a single bead, and the interaction parameters are chosen to
capture the relative strengths of canonical Watson-Crick base pairs (C-G and
A-U) as well as G-U wobble pairs. The total potential energy is written as

\begin{equation}
  U = U_{\mathrm{BA}} + U_{\mathrm{EV}} + U_{\mathrm{BP}},
  \label{eq:Utot}
\end{equation}

where $U_{\mathrm{BA}}$ denotes bonded backbone interactions,
$U_{\mathrm{EV}}$ denotes excluded-volume repulsion, and
$U_{\mathrm{BP}}$ denotes the many-body base-pairing interaction.

The solvent was modeled implicitly, so ionic effects were absorbed into the
effective coarse-grained interaction parameters rather than represented with
explicit ions.
% TODO: verify whether you want to state an effective salt concentration here.
% Reviewer specifically flagged the old 50 mM NaCl sentence.

%% --- Replace Bonded Interactions subsection with this ---

\subsection{Bonded Interactions}

The backbone covalent connectivity was modeled with a harmonic bond potential,

\begin{equation}
  U_{\mathrm{bond}} = k(r - r_0)^2,
\end{equation}

with spring constant $k = 7.5$~kcal\,mol$^{-1}$\,\AA$^{-2}$ and equilibrium
bond length $r_0 = 5.9$~\AA. Local backbone stiffness was modeled with a
harmonic angle potential,

\begin{equation}
  U_{\mathrm{angle}} = k_\theta(\theta - \theta_0)^2,
\end{equation}

with bending constant $k_\theta = 5.0$~kcal\,mol$^{-1}$\,rad$^{-2}$ and
equilibrium angle $\theta_0 = 150.0^\circ$, consistent with the extended
geometry of single-stranded nucleic-acid backbones.

Excluded-volume repulsion was represented using a shifted Lennard-Jones 12-6
potential with $\varepsilon = 2.0$~kcal\,mol$^{-1}$ and
$\sigma = 8.9090$~\AA. The potential was shifted to zero at a cutoff of
10.0~\AA\ using the \texttt{pair\_modify shift} command, retaining only the
repulsive core. The relatively large $\sigma$ reflects the physical size of a
nucleotide-level bead.

%% --- Replace Base-Pairing Potential subsection with this ---

\subsection{Base-Pairing Potential}

Base-pairing and stacking interactions were modeled using the custom LAMMPS
pair style \texttt{base/rna} as implemented in the SIS framework
used here~\cite{Aierken2024}, with a cutoff of 18.0~\AA. This interaction
represents the anisotropic geometry of RNA base pairing through inter-bead
distances, four angles, and two dihedral angles, thereby favoring A-form
helical geometry. Complementary C-G pairs were assigned
$\varepsilon = -5.0$~kcal\,mol$^{-1}$, whereas A-U and G-U wobble pairs were
assigned $\varepsilon = -3.33$~kcal\,mol$^{-1}$. All non-complementary pairs
were assigned $\varepsilon = 0$. The geometric parameters used for all base
pairs were an equilibrium contact distance of 3.0~\AA, an outer cutoff
distance of 13.8~\AA, and shape parameters
$(1.5,\;0.5,\;1.8326,\;0.9425,\;1.8326,\;1.1345)$.
% TODO: if your advisor wants this locked to the original model paper, re-check
% these numerical values against that source before final submission.

Non-bonded interactions between sites separated by one or two covalent links
(i.e., 1-2 and 1-3 pairs) were excluded to avoid double-counting bonded
interactions. All 1-4 non-bonded interactions were retained at full strength
using \texttt{special\_bonds lj 0.0 0.0 1.0}. Neighbor lists were built with a
binning algorithm using a skin distance of 15.0~\AA\ and were updated every
10 timesteps.

%% --- Replace NVT subsection with this ---

\subsection{NVT Equilibration}

Each minimized system was equilibrated for 50~ns
($5 \times 10^{6}$ timesteps with $\Delta t = 10$~fs) under constant-temperature
conditions. ASO simulations were run at 298~K or 300~K, depending on the
system.

Constant-temperature dynamics were generated using velocity-Verlet
integration together with a Langevin thermostat~\cite{Schneider1978}, as
implemented in LAMMPS. The Langevin damping constant was
$\tau = 100{,}000$~fs. Initial velocities were sampled from a Gaussian
distribution at the target temperature, and the linear momentum of the system
was zeroed every 1,000 steps.
% This wording is safer than implying simultaneous use of fix nvt and fix langevin.

%% --- Production Runs is mostly fine; keep your ref fix ---

\subsection{Production Runs}
\label{sec:production}

Sampling was performed using production runs for both RNA and ASO systems
under the same constant-temperature conditions, but over different time
scales: 1,000~ns for ASO simulations and 4~\textmu s for RNA simulations.
Restart files were written every $10^7$ steps, and the final configuration was
saved as a LAMMPS-formatted data file. Atomic positions were recorded every
50,000 timesteps (500~ps), including molecule IDs, atom IDs, atom types,
charges, and unwrapped Cartesian coordinates (\texttt{xu yu zu}). At least
five independent production runs were performed for each system used in the
binding-affinity analysis (Section~\ref{sec:binding_affinity}).

%% --- For simulated annealing, reviewer says it was not used in this paper.
%% Best fix: delete the subsection entirely unless it is genuinely part of the final workflow.